% GENERATED FILE. Edit canonical lessons, then run npm run generate:books.
% canonical-source-hash: fnv1a64:9f8bc325226afe49
% canonical-lessons: SA-C250-durgam, SA-C250-prasadah, SA-C250-kutiram, SA-C250-bhavanam, SA-C250-avasah

\chapter{Fort, Palace, Hut, Building, Lodging}
\label{ch:sa-fort-palace-hut-building-lodging}
\hlchaptermodality{250}

\begin{chapteropening}
\textbf{By the end of this chapter:} \emph{I can say a fort, a palace, a hut, a building and lodging.}

Complete the last lesson of chapter 250: I can say a fort, a palace, a hut, a building and lodging.
\end{chapteropening}

\section[\texorpdfstring{\sk{दुर्गम्} (durgam)}{दुर्गम् (durgam)}]{\sk{दुर्गम्} (durgam) — a fort}

\label{lesson:SA-C250-durgam}



\begin{quote}
Before the new one: say the Sanskrit for a film, a painted cloth, then the Sanskrit for a game.
\end{quote}

\subsection*{You'll want to know: \sk{दुर्गम्}}
\textbf{\sk{दुर्गम्}} — \emph{durgam} — \textquotedblleft{}a fort\textquotedblright{}.

\begin{grammarlens}[title={What you've built}]
One more word for this chapter.
\end{grammarlens}

\subsection*{Your turn}
\begin{itemize}
\raggedright
  \item \emph{Say it:} \emph{durgam}
  \item \emph{Say it:} \emph{durgam}, once more
  \item \emph{Say it:} say \emph{durgam}
  \item \emph{Recall:} say the Sanskrit for a film, a painted cloth, then the Sanskrit for a game, then say \emph{durgam} again
\end{itemize}

\begin{tcolorbox}[breakable,colback=teal!4,colframe=teal!35!black,title={Before you move on}]
What does \sk{दुर्गम्} mean? (\textquotedblleft{}A fort\textquotedblright{}.) Say it once more.
\end{tcolorbox}
\section[\texorpdfstring{\sk{प्रासादः} (prāsādaḥ)}{प्रासादः (prāsādaḥ)}]{\sk{प्रासादः} (prāsādaḥ) — a palace}

\label{lesson:SA-C250-prasadah}



\begin{quote}
Before the new one: say the Sanskrit for a game, then the Sanskrit for a fort.
\end{quote}

\subsection*{You'll want to know: \sk{प्रासादः}}
\textbf{\sk{प्रासादः}} — \emph{prāsādaḥ} — \textquotedblleft{}a palace\textquotedblright{}.

\begin{grammarlens}[title={What you've built}]
One more word for this chapter.
\end{grammarlens}

\subsection*{Your turn}
\begin{itemize}
\raggedright
  \item \emph{Say it:} \emph{prāsādaḥ}
  \item \emph{Say it:} \emph{prāsādaḥ}, once more
  \item \emph{Say it:} say \emph{prāsādaḥ}
  \item \emph{Recall:} say the Sanskrit for a game, then the Sanskrit for a fort, then say \emph{prāsādaḥ} again
\end{itemize}

\begin{tcolorbox}[breakable,colback=teal!4,colframe=teal!35!black,title={Before you move on}]
What does \sk{प्रासादः} mean? (\textquotedblleft{}A palace\textquotedblright{}.) Say it once more.
\end{tcolorbox}
\section[\texorpdfstring{\sk{कुटीरम्} (kuṭīram)}{कुटीरम् (kuṭīram)}]{\sk{कुटीरम्} (kuṭīram) — a hut, a cottage}

\label{lesson:SA-C250-kutiram}



\begin{quote}
Before the new one: say the Sanskrit for a fort, then the Sanskrit for a palace.
\end{quote}

\subsection*{You'll want to know: \sk{कुटीरम्}}
\textbf{\sk{कुटीरम्}} — \emph{kuṭīram} — \textquotedblleft{}a hut, a cottage\textquotedblright{}.

\begin{grammarlens}[title={What you've built}]
One more word for this chapter.
\end{grammarlens}

\subsection*{Your turn}
\begin{itemize}
\raggedright
  \item \emph{Say it:} \emph{kuṭīram}
  \item \emph{Say it:} \emph{kuṭīram}, once more
  \item \emph{Say it:} say \emph{kuṭīram}
  \item \emph{Recall:} say the Sanskrit for a fort, then the Sanskrit for a palace, then say \emph{kuṭīram} again
\end{itemize}

\begin{tcolorbox}[breakable,colback=teal!4,colframe=teal!35!black,title={Before you move on}]
What does \sk{कुटीरम्} mean? (\textquotedblleft{}A hut, a cottage\textquotedblright{}.) Say it once more.
\end{tcolorbox}
\section[\texorpdfstring{\sk{भवनम्} (bhavanam)}{भवनम् (bhavanam)}]{\sk{भवनम्} (bhavanam) — a building, a house}

\label{lesson:SA-C250-bhavanam}



\begin{quote}
Before the new one: say the Sanskrit for a palace, then the Sanskrit for a hut.
\end{quote}

\subsection*{You'll want to know: \sk{भवनम्}}
\textbf{\sk{भवनम्}} — \emph{bhavanam} — \textquotedblleft{}a building, a house\textquotedblright{}.

\begin{grammarlens}[title={What you've built}]
One more word for this chapter.
\end{grammarlens}

\subsection*{Your turn}
\begin{itemize}
\raggedright
  \item \emph{Say it:} \emph{bhavanam}
  \item \emph{Say it:} \emph{bhavanam}, once more
  \item \emph{Say it:} say \emph{bhavanam}
  \item \emph{Recall:} say the Sanskrit for a palace, then the Sanskrit for a hut, then say \emph{bhavanam} again
\end{itemize}

\begin{tcolorbox}[breakable,colback=teal!4,colframe=teal!35!black,title={Before you move on}]
What does \sk{भवनम्} mean? (\textquotedblleft{}A building, a house\textquotedblright{}.) Say it once more.
\end{tcolorbox}
\section[\texorpdfstring{\sk{आवासः} (āvāsaḥ)}{आवासः (āvāsaḥ)}]{\sk{आवासः} (āvāsaḥ) — lodging, a residence}

\label{lesson:SA-C250-avasah}



\begin{quote}
Before the new one: say the Sanskrit for a hut, then the Sanskrit for a building.
\end{quote}

\subsection*{You'll want to know: \sk{आवासः}}
\textbf{\sk{आवासः}} — \emph{āvāsaḥ} — \textquotedblleft{}lodging, a residence\textquotedblright{}.

\begin{grammarlens}[title={What you've built}]
One more word for this chapter.
\end{grammarlens}

\subsection*{Your turn}
\begin{itemize}
\raggedright
  \item \emph{Say it:} \emph{āvāsaḥ}
  \item \emph{Say it:} \emph{āvāsaḥ}, once more
  \item \emph{Say it:} say \emph{āvāsaḥ}
  \item \emph{Recall:} say the Sanskrit for a hut, then the Sanskrit for a building, then say \emph{āvāsaḥ} again
\end{itemize}

\begin{tcolorbox}[breakable,colback=teal!4,colframe=teal!35!black,title={Before you move on}]
What does \sk{आवासः} mean? (\textquotedblleft{}Lodging, a residence\textquotedblright{}.) Say it once more.
\end{tcolorbox}
